% Options for packages loaded elsewhere
\PassOptionsToPackage{unicode}{hyperref}
\PassOptionsToPackage{hyphens}{url}
\documentclass[
  a4paper,11pt]{article}
\usepackage{xcolor}
\usepackage[margin=1in]{geometry}
\usepackage{amsmath,amssymb}
\usepackage{graphicx}
\setcounter{secnumdepth}{-\maxdimen} % remove section numbering
\usepackage{iftex}
\ifPDFTeX
  \usepackage[T1]{fontenc}
  \usepackage[utf8]{inputenc}
  \usepackage{textcomp}
\else
  \usepackage{unicode-math}
  \defaultfontfeatures{Scale=MatchLowercase}
  \defaultfontfeatures[\rmfamily]{Ligatures=TeX,Scale=1}
\fi
\usepackage{lmodern}
\ifPDFTeX\else
  \setmainfont[]{Times New Roman}
\fi
\IfFileExists{upquote.sty}{\usepackage{upquote}}{}
\IfFileExists{microtype.sty}{%
  \usepackage[]{microtype}
  \UseMicrotypeSet[protrusion]{basicmath}
}{}
\makeatletter
\@ifundefined{KOMAClassName}{%
  \IfFileExists{parskip.sty}{%
    \usepackage{parskip}
  }{%
    \setlength{\parindent}{0pt}
    \setlength{\parskip}{6pt plus 2pt minus 1pt}}
}{% if KOMA class
  \KOMAoptions{parskip=half}}
\makeatother
\usepackage{longtable,booktabs,array}
\newcounter{none}
\usepackage{calc}
\usepackage{etoolbox}
\makeatletter
\patchcmd\longtable{\par}{\if@noskipsec\mbox{}\fi\par}{}{}
\makeatother
\IfFileExists{footnotehyper.sty}{\usepackage{footnotehyper}}{\usepackage{footnote}}
\makesavenoteenv{longtable}
\usepackage{bookmark}
\IfFileExists{xurl.sty}{\usepackage{xurl}}{}
\urlstyle{same}
\usepackage{hyperref}
\hypersetup{
  hidelinks,
  pdfcreator={LaTeX via pandoc}}
\setlength{\emergencystretch}{3em}
\providecommand{\tightlist}{%
  \setlength{\itemsep}{0pt}\setlength{\parskip}{0pt}}
\usepackage{amsthm}
\newtheorem{theorem}{Theorem}[section]
\newtheorem{proposition}[theorem]{Proposition}
\newtheorem{lemma}[theorem]{Lemma}
\newtheorem{corollary}[theorem]{Corollary}
\theoremstyle{definition}
\newtheorem{definition}[theorem]{Definition}
\theoremstyle{remark}
\newtheorem{remark}[theorem]{Remark}
\newtheorem{observation}[theorem]{Observation}
\begin{document}
\section{Shape-Invariant Standing Waves on Closed Spheres ---
Dimensional Stability and the Geometric Necessity of 3+1
Spacetime}\label{shape-invariant-standing-waves-on-closed-spheres-dimensional-stability-and-the-geometric-necessity-of-31-spacetime}

\textbf{Author}: Noriaki Kihara\\
\textbf{Affiliation}: WF System Co., Ltd.\\
\textbf{ORCID}: 0009-0004-6753-4020\\
\textbf{DOI}: https://doi.org/10.5281/zenodo.19731594\\
\textbf{Version}: v1.0\\
\textbf{Date}: April 2026

\begin{center}\rule{0.5\linewidth}{0.5pt}\end{center}

\subsection{Abstract}\label{abstract}

We systematically analyze the phase wave equation on closed
\(n\)-dimensional spheres \(S^n\) from one to four dimensions. We show
that the phase wave equation is independent of the spatial curvature
\(R\), and that \(R\) enters only through the standing wave condition
(the global boundary condition). We demonstrate that the parameter
description \((\theta_0, \lambda_{\text{base}}, A, N)\) of
shape-invariant standing waves is preserved across all dimensions, that
the conserved quantity \(\int |u|^2 \, dV\) carries the identity of the
wave packet, and that Huygens' principle holds only on odd-dimensional
spheres and fails on even-dimensional ones. From these facts, we show
that \textbf{the minimal nontrivial sphere on which shape-invariant
standing waves can stably exist is \(S^3\) (three-dimensional)}.
Furthermore, we add an independent argument from the discrete-algebraic
side via Frobenius' theorem and argue that both the continuous and
discrete sides necessitate a \textbf{3+1 spacetime structure}.

\textbf{Keywords}: phase wave equation, shape-invariant standing wave,
closed sphere, Huygens' principle, Frobenius' theorem, 3+1 spacetime

\begin{center}\rule{0.5\linewidth}{0.5pt}\end{center}

\subsection{§1 Introduction and Scope of
Claims}\label{introduction-and-scope-of-claims}

This paper describes, as purely geometric and number-theoretic facts,
the dimensional dependence of standing waves of the phase wave equation
on closed \(n\)-dimensional spheres \(S^n\).

The claims of this paper are limited to the following three points:

\begin{enumerate}
\def\labelenumi{\arabic{enumi}.}
\tightlist
\item
  The phase wave equation is an operation essentially orthogonal to the
  spatial curvature \(R\), and \(R\) enters only through the standing
  wave condition (the commensurability of the circumference of the
  closed space with the wavelength).
\item
  Shape-invariant standing waves that stably preserve their identity as
  wave packets exist only on odd-dimensional spheres where Huygens'
  principle holds. The minimal nontrivial dimension is \(S^3\).
\item
  By Frobenius' theorem (associative division algebras are limited to
  \(\mathbb{R}, \mathbb{C}, \mathbb{H}\)), four dimensions are
  independently necessitated from the discrete-algebraic side, and the
  \(1+3\) decomposition of the quaternions algebraically embeds the
  \(3+1\) spacetime structure.
\end{enumerate}

The following elements are \textbf{outside the scope of this paper}:

\begin{itemize}
\tightlist
\item
  Dynamical stability (eigenvalue analysis under time evolution,
  convexity of the energy functional)
\item
  Rigorous formulation as soliton solutions of nonlinear wave equations
\item
  Concrete correspondences with physical spacetime
\end{itemize}

The object called ``shape-invariant standing wave'' in this paper is
distinguished from ``soliton,'' which carries dynamical implications.
``Shape-invariant'' refers solely to the geometric property that the
amplitude distribution \(|\psi|\) is preserved under parallel transport,
and ``standing'' refers solely to the property that the nodal structure
is compatible with the arrangement of the family of spheres.

\begin{center}\rule{0.5\linewidth}{0.5pt}\end{center}

\subsection{\texorpdfstring{§2 Shape-Invariant Standing Waves on a
Closed String
(\(S^1\))}{§2 Shape-Invariant Standing Waves on a Closed String (S\^{}1)}}\label{shape-invariant-standing-waves-on-a-closed-string-s1}

\subsubsection{§2.1 Parameter Description}\label{parameter-description}

On a closed string (\(S^1\)) of circumference \(L\), a shape-invariant
standing wave is completely described by the following four parameters:

\begin{itemize}
\tightlist
\item
  \(\theta_0\): positional phase of the maximum wave height
\item
  \(\lambda_{\text{base}}\): fundamental mode wavelength
\item
  \(A\): amplitude of the maximum wave height
\item
  \(N\): number of modes in the Fourier expansion (truncation order)
\end{itemize}

Here \(N\) is the parameter controlling the ``roundness/sharpness'' of
the waveform:

\begin{itemize}
\tightlist
\item
  \(N = 1\): fundamental mode only (pure sinusoidal wave)
\item
  \(N\) intermediate: trapezoidal waveform with added odd-order
  harmonics
\item
  \(N \to \infty\): perfect square wave (fully localized pulse)
\end{itemize}

\subsubsection{§2.2 Standing Wave
Condition}\label{standing-wave-condition}

The necessary and sufficient condition for a standing wave to exist is:

\[L = k \cdot \lambda_{\text{base}} \quad (k \in \mathbb{Z}_{>0})\]

That is, the circumference of the closed string must be a positive
integer multiple of the fundamental mode wavelength. Under this
condition, the permitted wavelengths are automatically discretized to
\(\lambda_n = L/n\) (\(n = 1, 2, 3, \ldots\)).

\subsubsection{§2.3 Unification of Boundary
Conditions}\label{unification-of-boundary-conditions}

The phase wave equation requires the string to be closed but does not
specify the manner of closure. The difference between free-end
reflection (phase shift \(0\)) and fixed-end reflection (phase shift
\(\pi\)) can be treated uniformly, via the method of images, as the
even-parity/odd-parity subspaces of a closed string of length \(2L\).
The difference in boundary conditions is absorbed into the phase
selection (\(0\) or \(\pi\)) of the fundamental mode.

\begin{center}\rule{0.5\linewidth}{0.5pt}\end{center}

\subsection{\texorpdfstring{§3 Orthogonality of the Phase Wave Equation
and Curvature
\(R\)}{§3 Orthogonality of the Phase Wave Equation and Curvature R}}\label{orthogonality-of-the-phase-wave-equation-and-curvature-r}

\subsubsection{\texorpdfstring{§3.1
\(R\)-Independence}{§3.1 R-Independence}}\label{r-independence}

The phase wave equation is a local differential equation, and the
spatial curvature \(R\) (or equivalently \(L = 2\pi R\)) appears nowhere
in it. \(R\) enters only at \textbf{the level of the global boundary
condition}, which selects which wavelengths can constitute standing
waves:

\begin{itemize}
\tightlist
\item
  \(\lambda = L/n\) (\(n\) integer): phase shifts by \(2\pi n\) over one
  circuit, and a standing wave is established
\item
  \(\lambda > L\): phase shifts by less than \(2\pi\) over one circuit,
  and no standing wave forms
\item
  \(\lambda\) incommensurate with \(L\): phase never returns to its
  original value regardless of the number of circuits, and the
  trajectory fills the phase space quasi-periodically
\end{itemize}

In every case the wave equation itself is identical; whether \(R\) is
small or large, the wave obeys the same equation.

\subsubsection{\texorpdfstring{§3.2 A Posteriori Definition of
\(R\)}{§3.2 A Posteriori Definition of R}}\label{a-posteriori-definition-of-r}

From this structure, \(R\) is not a quantity within the wave equation
but rather \textbf{a geometric quantity that can be read off a
posteriori by surveying the set of standing-wave solutions}:

\[R = \frac{\lambda_{\text{base}}}{2\pi}\]

That is, \(R\) is ``\((2\pi)^{-1}\) times the longest wavelength
admitted as a standing wave'' --- a derived quantity, not an a priori
parameter.

\subsubsection{§3.3 Structural
Orthogonality}\label{structural-orthogonality}

\textbf{This is not a Euclidean approximation.} The phase wave equation
and the spatial curvature are essentially orthogonal operations, and
calculations in Euclidean coordinates are exactly correct. Curvature
\(R\) enters only through the global boundary condition (the
commensurability of the circumference of the closed space with the
wavelength) and has no effect whatsoever on the local solutions of the
wave equation.

\begin{center}\rule{0.5\linewidth}{0.5pt}\end{center}

\subsection{\texorpdfstring{§4 Extension to the Sphere \(S^2\) and
Antipodal
Convergence}{§4 Extension to the Sphere S\^{}2 and Antipodal Convergence}}\label{extension-to-the-sphere-s2-and-antipodal-convergence}

\subsubsection{§4.1 Preservation of Parameter
Structure}\label{preservation-of-parameter-structure}

Shape-invariant standing waves on the sphere \(S^2\) have essentially
the same parameter structure as on \(S^1\). By the \(\mathrm{SO}(3)\)
symmetry of \(S^2\), the standing wave condition takes the same form
regardless of which great circle is chosen:

\[2\pi R = n \cdot \lambda_{\text{base}}\]

The parameters change only in the dimension of the positional phase and
are exhausted by the four quantities
\((\theta_0, \lambda_{\text{base}}, A, N)\). The position of the wave
packet is a completely free parameter and is indistinguishable from
internal observation.

\subsubsection{§4.2 Antipodal Convergence
Phenomenon}\label{antipodal-convergence-phenomenon}

On \(S^2\), when the great-circle circumference is not an integer
multiple of \(\lambda\), a wave packet localized at a single point
executes the following motion:

\begin{enumerate}
\def\labelenumi{\arabic{enumi}.}
\tightlist
\item
  Initially: a peak localized at a single point
\item
  Expands as a ring-shaped (\(S^1\)) wavefront
\item
  Ring reaches maximum near the equator
\item
  Contracts toward the antipodal point on the opposite side
\item
  Converges again to a single point at the antipodal point
\item
  Repeats
\end{enumerate}

\subsubsection{§4.3 Geometric Dispersion}\label{geometric-dispersion}

The eigenfrequencies of \(S^2\) are \(\omega_l = c\sqrt{l(l+1)}/R\),
which are \textbf{non-equally spaced}. As a result, the wave packet is
strictly quasi-periodic rather than periodic, and the shape of the wave
packet deteriorates with each successive cycle.

\begin{center}\rule{0.5\linewidth}{0.5pt}\end{center}

\subsection{\texorpdfstring{§5 Conserved Quantity
\(\int |u|^2 \, dV\)}{§5 Conserved Quantity \textbackslash int \textbar u\textbar\^{}2 \textbackslash, dV}}\label{conserved-quantity-int-u2-dv}

\subsubsection{§5.1 Distinction Between Instantaneous Wave Height and
Conserved
Quantity}\label{distinction-between-instantaneous-wave-height-and-conserved-quantity}

In a perfectly lossless linear wave equation, the instantaneous wave
height \(u(x,t)\) itself is not a ``stable quantity'' either locally or
globally. What is conserved is the following nonlocal integral quantity:

\[\int_{S^n} |u|^2 \, dV = \sum_{l,m} |a_{lm}|^2 = \text{const}\]

The ``identity'' of a wave packet resides not in its instantaneous shape
but in this squared integral over the entire space.

\subsubsection{§5.2 Structural Isomorphism with the Schrödinger
Equation}\label{structural-isomorphism-with-the-schruxf6dinger-equation}

This conservation law is structurally isomorphic to
\(\int |\psi|^2 \, dV = 1\) in the Schrödinger equation of quantum
mechanics:

{\def\LTcaptype{none} % do not increment counter
\begin{longtable}[]{@{}
  >{\raggedright\arraybackslash}p{(\linewidth - 4\tabcolsep) * \real{0.3333}}
  >{\raggedright\arraybackslash}p{(\linewidth - 4\tabcolsep) * \real{0.3333}}
  >{\raggedright\arraybackslash}p{(\linewidth - 4\tabcolsep) * \real{0.3333}}@{}}
\toprule\noalign{}
\begin{minipage}[b]{\linewidth}\raggedright
\end{minipage} & \begin{minipage}[b]{\linewidth}\raggedright
Shape-invariant standing wave on a closed sphere
\end{minipage} & \begin{minipage}[b]{\linewidth}\raggedright
Schrödinger equation
\end{minipage} \\
\midrule\noalign{}
\endhead
\bottomrule\noalign{}
\endlastfoot
Field & \(u(x,t)\) (real-valued wave height) & \(\psi(x,t)\) (complex
amplitude) \\
Conserved density & \(u^2\) & \(|\psi|^2\) (probability density) \\
Conserved quantity & \(\int u^2 \, dA\) & \(\int |\psi|^2 \, dV = 1\) \\
Expansion basis & Spherical harmonics \(Y_{lm}\) & Hamiltonian
eigenfunctions \\
\end{longtable}
}

This isomorphism is not coincidental but a property arising from the
common mathematical skeleton of ``a linear equation expandable in an
orthogonal basis''; the distinction between classical and quantum is not
essential.

\begin{center}\rule{0.5\linewidth}{0.5pt}\end{center}

\subsection{\texorpdfstring{§6 Validity of Huygens' Principle on
\(S^3\)}{§6 Validity of Huygens' Principle on S\^{}3}}\label{validity-of-huygens-principle-on-s3}

\subsubsection{§6.1 Qualitative Difference Between Odd and Even
Dimensions}\label{qualitative-difference-between-odd-and-even-dimensions}

The propagation properties of the wave equation differ essentially
between even and odd spatial dimensions (Hadamard's result):

\begin{itemize}
\tightlist
\item
  \textbf{Odd dimensions} (\(S^1, S^3, S^5, \ldots\)): strict Huygens'
  principle holds. A point-source pulse propagates solely as a sharp
  wavefront, leaving no trailing wake.
\item
  \textbf{Even dimensions} (\(S^2, S^4, \ldots\)): Huygens' principle
  fails. A ``wake'' persists after the wavefront has passed.
\end{itemize}

\subsubsection{\texorpdfstring{§6.2 Shape-Invariant Standing Waves on
\(S^3\)}{§6.2 Shape-Invariant Standing Waves on S\^{}3}}\label{shape-invariant-standing-waves-on-s3}

On \(S^3\) (the hyperspherical shell of a perfectly symmetric
hypersphere in four-dimensional space):

\begin{itemize}
\tightlist
\item
  Conserved quantity: \(\int |u|^2 \, dV\) (\(V = 2\pi^2 R^3\))
\item
  Standing wave condition: \(2\pi R = n \cdot \lambda_{\text{base}}\)
  (\(\mathrm{SO}(4)\) symmetry)
\item
  Parameters:
  \((\theta_0, \phi_0, \psi_0, \lambda_{\text{base}}, A, N)\)
\end{itemize}

By Huygens' principle, a wave packet localized at a single point
\textbf{propagates solely as a sharp spherical shell (\(S^2\)-shaped
wavefront)}, leaving nothing behind. It converges exactly to a single
point at the antipodal point, where the peak is restored, and then
expands again as a spherical shell. Compared to the quasi-periodic
motion with afterglow on \(S^2\), a \textbf{cleaner restoration} is
achieved on \(S^3\).

\subsubsection{\texorpdfstring{§6.3 Failure of Huygens' Principle on
\(S^4\)}{§6.3 Failure of Huygens' Principle on S\^{}4}}\label{failure-of-huygens-principle-on-s4}

On \(S^4\), Huygens' principle fails again, and afterglow persists
behind the wavefront, as on \(S^2\). The clean restoration achieved on
\(S^3\) is lost.

\subsubsection{§6.4 Alternating Pattern of Dimension and
Stability}\label{alternating-pattern-of-dimension-and-stability}

{\def\LTcaptype{none} % do not increment counter
\begin{longtable}[]{@{}
  >{\raggedright\arraybackslash}p{(\linewidth - 4\tabcolsep) * \real{0.1579}}
  >{\raggedright\arraybackslash}p{(\linewidth - 4\tabcolsep) * \real{0.4211}}
  >{\raggedright\arraybackslash}p{(\linewidth - 4\tabcolsep) * \real{0.4211}}@{}}
\toprule\noalign{}
\begin{minipage}[b]{\linewidth}\raggedright
Dimension
\end{minipage} & \begin{minipage}[b]{\linewidth}\raggedright
Huygens' principle
\end{minipage} & \begin{minipage}[b]{\linewidth}\raggedright
Behavior of wave packet
\end{minipage} \\
\midrule\noalign{}
\endhead
\bottomrule\noalign{}
\endlastfoot
\(S^1\) & (one-dimensional) exact & Phase rotation only, perfectly
periodic \\
\(S^2\) & Fails & Antipodal convergence + afterglow, quasi-periodic \\
\(S^3\) & \textbf{Holds} & Sharp spherical-shell propagation, clean
restoration \\
\(S^4\) & Fails & Higher-dimensional analogue of \(S^2\), afterglow
present \\
\(S^5\) & \textbf{Holds} & Higher-dimensional analogue of \(S^3\),
clean \\
\end{longtable}
}

\begin{center}\rule{0.5\linewidth}{0.5pt}\end{center}

\subsection{§7 Conditions for Stable Existence of Shape-Invariant
Standing
Waves}\label{conditions-for-stable-existence-of-shape-invariant-standing-waves}

\subsubsection{§7.1 Loss of Wave Packet Identity in Even
Dimensions}\label{loss-of-wave-packet-identity-in-even-dimensions}

We organize the phenomena occurring in even dimensions (\(S^2, S^4\)):

\begin{itemize}
\tightlist
\item
  The conserved quantity \(\int |u|^2 \, dV\) is exactly conserved
\item
  Each individual mode amplitude \(|a_{lm}|\) is also exactly conserved
\item
  However, \textbf{the identity as a localized wave packet} is not
  maintained
\end{itemize}

Because the eigenfrequencies \(\omega_l\) are arranged in irrational
ratios, the phases of the individual modes drift independently,
undergoing quasi-periodic motion in phase space. Under long-time
averaging, a distribution spread over the entire sphere becomes the
typical state, with localization partially reviving by occasional
coincidental phase alignment (Poincaré recurrence). This is consistent
with the picture of a quasi-periodic dynamical system in the sense of
KAM theory.

\subsubsection{§7.2 Necessary and Sufficient Conditions for Stable
Shape-Invariant Standing
Waves}\label{necessary-and-sufficient-conditions-for-stable-shape-invariant-standing-waves}

The conditions for a shape-invariant standing wave to ``stably exist
while preserving its identity as a wave packet'' are:

\begin{enumerate}
\def\labelenumi{\arabic{enumi}.}
\tightlist
\item
  \textbf{Huygens' principle holds}: the wavefront leaves no trailing
  wake
\item
  \textbf{Eigenfrequencies are close to integer ratios}: phase coherence
  is regenerated
\end{enumerate}

These are simultaneously satisfied \textbf{only on odd-dimensional
spheres}.

\subsubsection{§7.3 Minimality}\label{minimality}

\(S^1\) trivially satisfies Huygens' principle, but in one dimension it
exhibits only the degenerate behavior of phase rotation. The minimal
dimension that supports nontrivial wave packet formation, restoration,
and antipodal convergence is \textbf{\(S^3\)}.

\textbf{Proposition}: The minimal dimension \(n\) on which
shape-invariant standing waves can exist nontrivially and stably on a
closed \(n\)-dimensional sphere is \(n = 3\).

\begin{center}\rule{0.5\linewidth}{0.5pt}\end{center}

\subsection{§8 Independent Argument via Frobenius'
Theorem}\label{independent-argument-via-frobenius-theorem}

\subsubsection{§8.1 Associative Division
Algebras}\label{associative-division-algebras}

By Frobenius' theorem (1877), the finite-dimensional associative
division algebras over the real numbers are limited to the following
three:

\begin{itemize}
\tightlist
\item
  \(\mathbb{R}\) (real numbers, 1-dimensional)
\item
  \(\mathbb{C}\) (complex numbers, 2-dimensional)
\item
  \(\mathbb{H}\) (quaternions, 4-dimensional)
\end{itemize}

No natural division algebra exists in three dimensions. Extending while
preserving algebraic structure (addition, subtraction, multiplication,
division, and norm) skips three dimensions and arrives at four.

\subsubsection{\texorpdfstring{§8.2 The \(1+3\) Decomposition of
Quaternions}{§8.2 The 1+3 Decomposition of Quaternions}}\label{the-13-decomposition-of-quaternions}

A quaternion \(q = a + bi + cj + dk\) decomposes intrinsically within
the algebraic structure into \(1\) dimension (scalar part \(a\)) \(+ 3\)
dimensions (vector part \(bi + cj + dk\)). This decomposition is not
arbitrary but a natural structure demanded by the conjugation, norm, and
multiplication of quaternions.

\subsubsection{§8.3 Upper Bound: Non-Associativity of
Octonions}\label{upper-bound-non-associativity-of-octonions}

Extending to the 8-dimensional octonions \(\mathbb{O}\) breaks the
associative law (Hurwitz's theorem). Since losing associativity makes it
difficult to describe the composition and interaction of waves, four
dimensions constitute the maximal dimension for an associative division
algebra.

\begin{center}\rule{0.5\linewidth}{0.5pt}\end{center}

\subsection{§9 Convergence of the Dual
Argument}\label{convergence-of-the-dual-argument}

\subsubsection{§9.1 Independent Arguments from the Continuous and
Discrete
Sides}\label{independent-arguments-from-the-continuous-and-discrete-sides}

{\def\LTcaptype{none} % do not increment counter
\begin{longtable}[]{@{}
  >{\raggedright\arraybackslash}p{(\linewidth - 4\tabcolsep) * \real{0.3333}}
  >{\raggedright\arraybackslash}p{(\linewidth - 4\tabcolsep) * \real{0.3333}}
  >{\raggedright\arraybackslash}p{(\linewidth - 4\tabcolsep) * \real{0.3333}}@{}}
\toprule\noalign{}
\begin{minipage}[b]{\linewidth}\raggedright
Aspect
\end{minipage} & \begin{minipage}[b]{\linewidth}\raggedright
Conclusion
\end{minipage} & \begin{minipage}[b]{\linewidth}\raggedright
Basis
\end{minipage} \\
\midrule\noalign{}
\endhead
\bottomrule\noalign{}
\endlastfoot
Continuous waves (phase wave equation on \(S^n\)) & Three-dimensional
space is preferred & Huygens' principle in odd dimensions \\
Discrete algebra (lattice structure) & Four dimensions are necessitated
& Frobenius' theorem \\
Synthesis & \textbf{3 + 1 = 4} & \(1+3\) decomposition of quaternions \\
\end{longtable}
}

\subsubsection{§9.2 Unified Proposition}\label{unified-proposition}

The ``shape-invariant standing waves stable in odd dimensions'' on the
continuous side and ``Frobenius' theorem'' on the discrete side, though
seemingly independent, are two facets of a single requirement:
\textbf{the minimal spacetime in which stable shape-invariant standing
waves can exist}:

\begin{enumerate}
\def\labelenumi{\arabic{enumi}.}
\tightlist
\item
  \textbf{Spatial structure}: a space that geometrically protects the
  identity of wave packets → odd-dimensional spheres → the minimal
  nontrivial dimension is \(S^3\) (3-dimensional)
\item
  \textbf{Algebraic structure}: interactions of wave packets close
  within the same category → associative division algebras → the minimal
  nontrivial dimension is \(\mathbb{H}\) (4-dimensional)
\item
  \textbf{Synthesis}: via the \(1+3\) decomposition of quaternions,
  unified as 3-dimensional space \(+\) 1-dimensional time axis
\end{enumerate}

\textbf{Theorem (Convergence of the dual argument)}: The minimal closed
space on which shape-invariant standing waves stably exist and possess
algebraically closed interactions is \(S^3\), and the minimal structure
describing this within an associative division algebra is the \(1+3\)
decomposition of the quaternions \(\mathbb{H}\), namely the \(3+1\)
spacetime structure.

\begin{center}\rule{0.5\linewidth}{0.5pt}\end{center}

\subsection{§10 Orthogonality of the Phase Wave Equation and Spatial
Curvature}\label{orthogonality-of-the-phase-wave-equation-and-spatial-curvature}

The orthogonality between the phase wave equation and the curvature
\(R\), demonstrated in §3, is preserved in higher dimensions. For
\(S^1 \to S^2 \to S^3 \to S^4\) alike:

\begin{itemize}
\tightlist
\item
  The wave equation is a universal equation independent of \(R\)
\item
  \(R\) is a geometric quantity determined a posteriori through the
  standing wave condition
\item
  Calculations in Euclidean coordinates are exactly correct (not an
  approximation)
\end{itemize}

This orthogonality means that shape-invariant standing waves are
indifferent to spatial distortion. Waves propagate along geodesics, but
whether the space is distorted or not has no effect on the local
solutions of the wave equation.

\begin{center}\rule{0.5\linewidth}{0.5pt}\end{center}

\subsection{§11 Summary}\label{summary}

We summarize the main results of this paper:

\begin{enumerate}
\def\labelenumi{\arabic{enumi}.}
\item
  \textbf{\(R\)-independence}: The phase wave equation and the spatial
  curvature are essentially orthogonal, and \(R\) is a geometric
  quantity that can be read off a posteriori from the set of
  standing-wave solutions.
\item
  \textbf{Universality of the parameter description}: The parameters
  \((\theta_0, \lambda_{\text{base}}, A, N)\) of shape-invariant
  standing waves are described by the same skeleton from \(S^1\) to
  \(S^n\), with only the dimension of the positional phase varying.
\item
  \textbf{Conserved quantity}: The identity of a wave packet is carried
  not by the instantaneous wave height \(u\) but by the nonlocal squared
  integral \(\int |u|^2 \, dV\).
\item
  \textbf{Dimensional selection via Huygens' principle}: Shape-invariant
  standing waves can stably exist only on odd-dimensional spheres, and
  the minimal nontrivial dimension is \(S^3\) (three-dimensional).
\item
  \textbf{Convergence of the dual argument}: The continuous side
  (Huygens' principle) and the discrete side (Frobenius' theorem)
  independently necessitate the same \(3+1\) spacetime structure.
\end{enumerate}

\begin{center}\rule{0.5\linewidth}{0.5pt}\end{center}

\subsection{References}\label{references}

\begin{itemize}
\tightlist
\item
  {[}P1{]} Noriaki Kihara, ``Central Projection Between Hyperspheres,''
  2025.
\item
  {[}P2{]} Noriaki Kihara, ``Regularity of Central Projection Between
  Hyperspherical Shells,'' DOI: 10.5281/zenodo.19692192, 2026.
\item
  {[}W8{]} Noriaki Kihara, ``Six-Dimensional Hyperrectangular
  Parallelepiped and Particle Model,'' 2025.
\item
  Hadamard, J., \emph{Lectures on Cauchy's Problem in Linear Partial
  Differential Equations}, 1923.
\item
  Frobenius, G., ``Über lineare Substitutionen und bilineare Formen'',
  \emph{J. reine angew. Math.}, \textbf{84}, 1--63, 1877.
\item
  Hurwitz, A., ``Über die Composition der quadratischen Formen von
  beliebig vielen Variablen'', \emph{Nachr. Ges. Wiss. Göttingen},
  309--316, 1898.
\end{itemize}
\end{document}
